\section{Experimental Setup}
\label{sec:experimental-setup}

We design experiments to answer four questions. \textbf{Q1:} Does the representation-based factorization (ResponseVec) outperform direct-generation and strong-prompting baselines on within-instrument prediction, in both accuracy and calibration? \textbf{Q2:} How does the advantage degrade under the three transfer regimes---unseen questions, cross-domain items, and out-of-distribution demographic intersections---and does the degradation differ between ResponseVec and RespondentVec? \textbf{Q3:} How do the baseline families compare on the accuracy--cost Pareto plane, and where does each regime sit on the frontier? \textbf{Q4:} Which design choices (generative vs.\ discriminative encoder, option-aware vs.\ option-agnostic decoder, history, demographics, pooling) drive the observed performance?

\subsection{Data}
\label{sec:data}

We use the \textbf{SocioBench} benchmark built from the International Social Survey Programme (ISSP), which provides nationally representative survey responses across four substantively distinct domains---Environment, Role of Government, Social Inequality, and Work Orientations---spanning approximately 30 countries with roughly 500 respondents per domain. Each respondent has a verified demographic profile (sex, age, education, income, employment, marital status, urbanicity) and item-level ordinal answers with full option sets, satisfying the requirements of our formalization in Section~\ref{sec:problem}. SocioBench is publicly available, documented, and already integrated into our pipeline; we chose it over single-country panels (e.g., ANES) precisely because the multi-country, multi-domain structure supports both cross-domain and cross-demographic transfer within one dataset, avoiding the ingestion and re-interview-availability risk that would attend a panel-wave design. After applying the minimum-answered-items filter (at least 10 items per respondent) and capping at 20 items per domain, the working dataset comprises on the order of $2\times10^3$ respondents and $10^2$ items across the four domains.

\subsection{Three Leakage-Safe Transfer Regimes}
\label{sec:regimes}

The three transfer regimes are the backbone of our evaluation, and each is enforced by hard leakage assertions in the code. \emph{R1 Unseen questions (Protocol B):} within each domain, items are partitioned into a fixed calibration pool (40\% of items, capped at 8) plus six disjoint target folds; each candidate target item is a final-test item in exactly one outer fold, while the decoder trains on four folds and validates on one. History is drawn exclusively from the calibration pool, so no target-fold item can ever enter any history. This regime isolates item-level generalization while holding the respondent population fixed. \emph{R2 Cross-domain (Protocol C):} a leave-one-domain-out protocol in which the decoder trains on three domains and is tested on the held-out domain's seen items, with history drawn from the respondent's other answered items in that domain. This regime tests whether the representation and decoder transfer across substantively distinct content domains with different option sets. \emph{R3 OOD-demographic-intersection (Protocol D):} entire demographic-intersection cells (defined by sex $\times$ age-bin $\times$ education) are held out into a dedicated \texttt{ood\_intersection} split---training never sees any respondent from those cells---and the held-out respondents are scored on their seen items, with history drawn from their own other answers. This regime tests generalization to respondents whose demographic combinations were never observed at training, a setting that the silicon-sampling literature has not systematically evaluated and that complements the cross-survey transfer gap identified by \citet{ku2026silicon}. All three regimes share the same frozen item folds, demographic template, history-retrieval procedure, and option-order randomization, so differences stem from the prediction mechanism rather than from data-handling artifacts.

\subsection{Baselines}
\label{sec:baselines}

We compare against baseline families spanning the methodological space, all of which are implemented in our codebase and share the same prompt, frozen option encoder, and leakage-safe splits. \textbf{B1 Direct generation:} the LLM is prompted with $(D_r, H_r, q^\ast, o^\ast)$ and its option-token logits are read out; we report five progressively calibrated variants---raw softmax, temperature-scaled, prior-blended, scalar-calibrated, and a supervised option-aware neural calibrator---so that a ResponseVec win cannot be explained by unequal calibration access. \textbf{B2 Majority and MIRT:} the non-personalized per-question majority floor, plus a demographic-prior multidimensional item-response model refined by MAP steps on the $K$ known answers, the classical psychometric frontier. \textbf{B3 Item-conditional CF:} a leakage-safe collaborative-filtering model that tilts the prior with shrunk empirical conditionals $P(y_j\mid y_k)$, weighted by the shrunk item--item Spearman graph. \textbf{B4 QLoRA fine-tuning:} the 8B backbone is quantized to 4-bit and a rank-16 adapter is trained on the same canonical prompts with randomized option labels, representing the task-adaptivity upper bound at high compute. \textbf{B5 Generic text vector:} the canonical prompt is encoded by a frozen BGE embedder and scored by the same option-aware decoder as ResponseVec, isolating the contribution of the LLM-based encoder. \textbf{B6 Raw hidden states:} the last-layer and mean-pool hidden states of the LLM on the canonical prompt feed the option-aware decoder, isolating the contribution of contrastive alignment. \textbf{B7 LLM2Vec (no Gen):} the discriminative LLM2Vec encoder replaces LLM2Vec-Gen in our pipeline, isolating the contribution of preserving output semantics. Our methods \textbf{M1 ResponseVec} and \textbf{M2 RespondentVec} complete the comparison. We do not implement a separate generative-agent (GEMS) baseline: its memory-stream architecture collapses, under our canonical prompt, into the strong-prompting path already covered by B1 with history, and a faithful re-implementation would constitute a separate study; we discuss this scope choice honestly in Section~\ref{sec:discussion}.

\subsection{Metrics}
\label{sec:metrics}

We report \emph{accuracy} and \emph{normalized ordinal error} for choice prediction, \emph{NLL} and \emph{Ranked Probability Score (RPS)} as proper scoring rules for ordinal outcomes, and \emph{Brier score} for probabilistic prediction. For calibration we report \emph{Expected Calibration Error (ECE)} with 15-bin reliability curves. For transfer we report the accuracy drop $\Delta_{\text{transfer}} = \text{Acc}_{\text{R1}} - \text{Acc}_{\text{R}k}$ and the transfer ratio $\rho = \text{Acc}_{\text{R}k}/\text{Acc}_{\text{R1}}$ relative to the unseen-item regime. For robustness we report the variance of accuracy under option-permutation perturbations (three deterministic semantic permutations). For compute we report GPU-hours per 1000 respondents, wall-clock latency per respondent, and peak GPU memory, all measured on the same A100 40GB hardware. All accuracy differences are reported with paired two-way (respondent $\times$ item) cluster bootstrap 95\% confidence intervals and Holm-corrected family-wise $p$-values, and all methods are run over three decoder seeds and three option seeds with means and standard deviations reported.

\subsection{Implementation}
\label{sec:implementation}

The LLM backbone is Qwen3-8B; the generative embedder is LLM2Vec-Gen initialized from the same backbone, with a 0.6B small-backbone variant for the scale analysis. The decoder is a bilinear option-aware head with projection dimension 256, dropout 0.10, learned temperature, and a fixed prior-blend weight $\eta$ selected once on validation and shared across representation families; the option-agnostic ablation replaces the bilinear head with a two-layer MLP. The decoder is trained with Adam (lr $10^{-3}$, weight decay $10^{-5}$) for at most 100 epochs with early stopping (patience 10) on validation NLL, plus an ordinal RPS auxiliary loss ($\lambda=0.1$). Calibration temperature and per-option biases are fit on a held-out validation set. For B4 (QLoRA) we use 4-bit NF4 quantization with rank-16, $\alpha=32$ adapters on \texttt{q\_proj}/\texttt{v\_proj}, trained for 1200 steps at effective batch 64. Representation caches are fingerprinted by model revision, quantization, item partition, and candidate-row hashes, and are immutable and resumable. Code, configurations, and frozen model revisions will be released upon acceptance to support replication of the full protocol.
